\section{Planning prévisionnel}

\begin{table}[H]
\centering
\begin{tabular}{|c|c|p{9cm}|}
\hline
\textbf{Sem.} & \textbf{Heures} & \textbf{Tâches} \\
\hline
1  & 8h  & Setup poste, Git, outils. Rencontres superviseurs. Analyse infra. Brouillon cahier des charges. \\
\hline
2  & 8h  & Étude Ray, SAM3, alternatives S3, PostGIS. Base du cahier des charges. \\
\hline
3  & 8h  & Dockerfile SAM3. Demande accès HuggingFace. Premiers tests sur pod K8s GPU. \\
\hline
4  & 8h  & Build image Docker. Push registry. Inférence SAM3 validée sur cluster. \\
\hline
5  & 8h  & Prototype Ray distribué. Intégration Ray + SAM3. \\
\hline
6  & 8h  & Client S3 boto3. Connexion MinIO. Pipeline S3 $\rightarrow$ Ray $\rightarrow$ SAM3. \\
\hline
7  & 8h  & Scénario batch : Ray Data sur 2000 images. Stockage Parquet sur S3. \\
\hline
\textbf{8}  & \textbf{8h}  & \textbf{Scénario inférence à la demande. Dockerisation complète. Remise cahier des charges.} \\
\hline
9  & 8h  & Observabilité : Prometheus + Grafana sur K8s. Métriques GPU/mémoire/latence. \\
\hline
10 & 8h  & Rancher : configuration cluster, gestion namespaces, monitoring. \\
\hline
11 & 8h  & Intégration Label Studio : API, export annotations, tests validation humaine. \\
\hline
12 & 8h  & K8s manifests complets (Jobs Ray, Deployments, PVC). Tests intégration E2E. \\
\hline
13 & 8h  & Tests pipeline complet (2 scénarios). Corrections. Préparation rapport intermédiaire. \\
\hline
\textbf{14} & \textbf{8h}  & \textbf{Remise rapport intermédiaire. Feedback superviseurs.} \\
\hline
15 & 8h  & Intégration Label Studio. Export annotations. Ajustements post-feedback. \\
\hline
16 & 8h  & Optimisations performance. Benchmarks downsampling. \\
\hline
17 & 8h  & Documentation déploiement. Validation pipeline sur cluster HEIG. \\
\hline
\textbf{18} & \textbf{63h} & \textbf{TB à plein temps. Déploiement production. Premiers runs à grande échelle.} \\
\hline
19 & 63h & Tests GPU cluster. Benchmarks E2E (scalabilité, précision, ressources). Tuning. \\
\hline
20 & 63h & Rédaction rapport : Introduction, contexte, architecture, choix techniques. \\
\hline
21 & 63h & Rédaction rapport : Implémentation, résultats, benchmarks, conclusions. \\
\hline
\textbf{22} & \textbf{62h} & \textbf{Finalisation rapport. Résumé publiable. Relecture. Soumission GAPS.} \\
\hline
\multicolumn{2}{|c|}{\textbf{Total}} & \textbf{450 heures} \\
\hline
\end{tabular}
\caption{Planning prévisionnel --- 450 heures sur 22 semaines}
\label{tab:planning}
\end{table}

\section{Livrables}

\begin{enumerate}
    \item \textbf{Cahier des charges} : 09 avril 2026
    \item \textbf{Rapport intermédiaire} : 20 mai 2026
    \item \textbf{Code source} : Repository Git (Ray, SAM3, MinIO, Parquet, K8s, observabilité)
    \item \textbf{Rapport de Bachelor} : 23 juillet 2026
    \item \textbf{Résumé publiable} : 23 juillet 2026
    \item \textbf{Présentation orale} : août/septembre 2026
\end{enumerate}

\section{Jalons importants}

\begin{table}[H]
\centering
\begin{tabular}{|l|p{9cm}|}
\hline
\textbf{Date} & \textbf{État attendu du pipeline} \\
\hline
09 avril 2026 & SAM3 sur cluster GPU, Ray + S3 fonctionnels, scénarios batch et on-demand prototypés. \\
\hline
20 mai 2026   & Pipeline E2E complet (MinIO + Ray + SAM3 + Parquet S3 + Observabilité + Docker + K8s). \\
\hline
23 juillet 2026 & Pipeline déployé sur cluster HEIG avec GPU. Rapport final soumis. \\
\hline
Août--Sept 2026 & Soutenance avec démo live du pipeline. \\
\hline
\end{tabular}
\caption{Jalons officiels HEIG-VD}
\label{tab:jalons}
\end{table}
